\documentclass[aspectratio=169,8pt]{beamer}
\usetheme{Madrid}
\usepackage[utf8]{inputenc}
\usepackage[T1]{fontenc}
\usepackage{amsmath,amssymb}
\usepackage{booktabs}
\usepackage{enumitem}
\setlist[enumerate]{label=\Alph*.,leftmargin=*,itemsep=1pt,topsep=2pt}
\setlist[itemize]{leftmargin=*,itemsep=1pt,topsep=2pt}
\setbeamerfont{block body}{size=\tiny}
\setbeamerfont{block title}{size=\scriptsize}
\setbeamerfont{frametitle}{size=\footnotesize}
\setbeamerfont{normal text}{size=\tiny}
\setbeamertemplate{navigation symbols}{}
\setbeamertemplate{itemize item}{\tiny$\bullet$}
\setbeamertemplate{itemize subitem}{\tiny$\circ$}

\title[GARP RAI]{GARP Risk \& AI --- Q551--Q650}
\subtitle{Unofficial Exam Prep}
\author{Unofficial}
\date{\today}

\begin{document}
	
	\begin{frame}
		\titlepage
	\end{frame}
	
	% =================================================================
	\section{Advanced AI Risk Taxonomy}
	% =================================================================
	

	
	
	\begin{frame}{Q560 --- AI Risk Taxonomy Categories}
		\begin{block}{Question}
			Which is a common category in AI risk taxonomies?
			\begin{enumerate}
				\item Technical risk, encompassing adversarial robustness, distributional shift, data-provenance integrity, and reproducibility variance.
				\item Ethical risk, encompassing allocative harm, representational harm, transparency deficits, and human-oversight sufficiency.
				\item Operational risk, encompassing third-party concentration, model-drift telemetry, process-failure modes, and resource-exhaustion scenarios.
				\item All of the above are common categories in AI risk taxonomies, alongside legal, regulatory, and reputational dimensions.
			\end{enumerate}
		\end{block}
		\begin{exampleblock}{Answer: D}\end{exampleblock}
		\begin{block}{Explanation}
			AI risk taxonomies span technical, ethical, operational, and other categories; all are common.
		\end{block}
	\end{frame}
	
	% =================================================================
	\section{AI Model Risk in Financial Services}
	% =================================================================
	
	\begin{frame}{Q561 --- Credit Scoring Model Risk}
		\begin{block}{Question}
			Which is the most relevant model risk for an AI-based credit scoring system?
			\begin{enumerate}
				\item The colour scheme and layout of the user interface may not align with brand guidelines, leading to inconsistent customer experience and reduced engagement metrics.
				\item Potential unfairness, inaccuracy, and adverse regulatory consequences from biased or poorly performing predictions, including proxy discrimination, explainability gaps, and ECOA/FCRA non-compliance.
				\item The storage cost of retaining historical training data may exceed budgeted cloud infrastructure allocations over multi-year horizons, impacting FinOps unit economics.
				\item Marketing effectiveness may decline if the model's output is not integrated with campaign management platforms, customer segmentation tools, and CRM workflows.
			\end{enumerate}
		\end{block}
		\begin{exampleblock}{Answer: B}\end{exampleblock}
		\begin{block}{Explanation}
			Credit scoring model risk includes fairness, accuracy, and regulatory risk. Options A, C, and D are secondary or irrelevant.
		\end{block}
	\end{frame}
	
	\begin{frame}{Q562 --- Fraud Detection Model Risk}
		\begin{block}{Question}
			For an AI fraud detection model, which is the most critical model risk?
			\begin{enumerate}
				\item Overfitting to historical fraud patterns, missing novel fraud schemes, and failing to adapt to adversarial concept drift under non-stationary distributions.
				\item Storage cost of retaining transaction-level features and model artefacts across multiple retraining cycles, impacting data-lake tiering strategies.
				\item User interface design that may not clearly present fraud alerts to analysts, slowing investigation workflows and increasing mean-time-to-resolution.
				\item Training time may exceed operational windows, delaying model refresh cycles and reducing responsiveness to emerging fraud typologies.
			\end{enumerate}
		\end{block}
		\begin{exampleblock}{Answer: A}\end{exampleblock}
		\begin{block}{Explanation}
			Overfitting to past patterns can miss novel fraud. Options B, C, and D are less critical.
		\end{block}
	\end{frame}
	
	
	\begin{frame}{Q564 --- Market Risk Model}
		\begin{block}{Question}
			Which is a typical risk for an AI-enhanced market risk model?
			\begin{enumerate}
				\item Training on non-stationary data leading to unstable predictions and regime-change failures under structural breaks.
				\item Data storage costs may escalate if tick-level market data is retained for extended backtesting periods, impacting data-lake economics.
				\item Logo design may not comply with brand guidelines, creating inconsistent visual identity across risk reports and dashboards.
				\item Employee satisfaction may decline if model outputs are not integrated with existing workflow tools and trading-desk platforms.
			\end{enumerate}
		\end{block}
		\begin{exampleblock}{Answer: A}\end{exampleblock}
		\begin{block}{Explanation}
			Non-stationarity of financial markets is a key challenge. Options B, C, and D are unrelated.
		\end{block}
	\end{frame}
	
	
	\begin{frame}{Q566 --- Model Fairness in Lending}
		\begin{block}{Question}
			Which is a recognised fairness metric used in lending AI?
			\begin{enumerate}
				\item Demographic parity, measured as the ratio of positive prediction rates between protected and non-protected groups.
				\item Storage cost parity, measured as the ratio of data storage costs between protected and non-protected groups.
				\item Latency parity, measured as the ratio of inference times between protected and non-protected groups.
				\item Colour parity, measured as the ratio of UI colour contrast ratios between protected and non-protected groups.
			\end{enumerate}
		\end{block}
		\begin{exampleblock}{Answer: A}\end{exampleblock}
		\begin{block}{Explanation}
			Demographic parity is a standard fairness metric. Options B, C, and D are not.
		\end{block}
	\end{frame}
	
	\begin{frame}{Q567 --- Algorithmic Trading Model Risk}
		\begin{block}{Question}
			Which is a key risk of algorithmic trading models?
			\begin{enumerate}
				\item Flash-crash amplification, unintended feedback loops, and market impact from correlated strategies under stressed liquidity conditions.
				\item Storage capacity may be insufficient to retain tick-level data for regulatory reporting and backtesting, impacting MiFID II compliance.
				\item Logo design may not comply with exchange branding requirements, creating regulatory filing issues and venue-certification delays.
				\item Employee benefits may not be competitive, leading to talent attrition in quantitative teams and loss of institutional knowledge.
			\end{enumerate}
		\end{block}
		\begin{exampleblock}{Answer: A}\end{exampleblock}
		\begin{block}{Explanation}
			Algorithmic trading risks include flash crashes and correlated behaviour. Options B, C, and D are unrelated.
		\end{block}
	\end{frame}
	
	\begin{frame}{Q568 --- Chatbot Risk in Banking}
		\begin{block}{Question}
			Which is the most significant risk from deploying a GenAI chatbot in banking?
			\begin{enumerate}
				\item Generating incorrect or non-compliant advice to customers, creating legal and reputational risk under consumer-protection regimes.
				\item Faster response times may reduce customer wait times but could increase volume of unverified interactions and hallucination exposure.
				\item Cost savings may improve margins but could reduce investment in human oversight and quality assurance.
				\item Customer satisfaction may increase but could mask underlying model deficiencies and bias.
			\end{enumerate}
		\end{block}
		\begin{exampleblock}{Answer: A}\end{exampleblock}
		\begin{block}{Explanation}
			Advice quality and compliance are central concerns. Options B, C, and D are benefits.
		\end{block}
	\end{frame}
	
	\begin{frame}{Q569 --- Model Risk from Alternative Data}
		\begin{block}{Question}
			Which is a key governance concern when using alternative data in financial AI?
			\begin{enumerate}
				\item Legal right to use the data, privacy, and potential for discriminatory impact through proxy variables and disparate-impact mechanisms.
				\item Storage cost may exceed budgeted cloud infrastructure allocations over multi-year horizons, impacting FinOps unit economics.
				\item Model speed may decline if alternative data requires extensive preprocessing and feature engineering pipelines.
				\item Logo design may not comply with brand guidelines, creating inconsistent visual identity across risk reports.
			\end{enumerate}
		\end{block}
		\begin{exampleblock}{Answer: A}\end{exampleblock}
		\begin{block}{Explanation}
			Legal rights, privacy, and fairness are key governance concerns. Options B, C, and D are secondary.
		\end{block}
	\end{frame}
	
	
	% =================================================================
	\section{AI Model Lifecycle and Change Management}
	% =================================================================

	
	

	


	
	\begin{frame}{Q580 --- A/B Testing of Models}
		\begin{block}{Question}
			What is A/B testing in the AI deployment context?
			\begin{enumerate}
				\item Randomly assigning users to two model versions and comparing outcomes under controlled experimental conditions.
				\item Retiring both models, as A/B testing is a decommissioning technique and archival strategy.
				\item Doubling training data, as A/B testing is a data augmentation technique and synthetic-data generation method.
				\item Storing data twice, as A/B testing is a storage redundancy technique and disaster-recovery strategy.
			\end{enumerate}
		\end{block}
		\begin{exampleblock}{Answer: A}\end{exampleblock}
		\begin{block}{Explanation}
			A/B testing compares model versions via controlled experiments. Options B, C, and D are unrelated.
		\end{block}
	\end{frame}
	
	% =================================================================
	\section{AI Incident Management}
	% =================================================================

	
	\begin{frame}{Q582 --- AI Incident Root Cause}
		\begin{block}{Question}
			An AI model begins producing biased outputs. Which is the most likely root cause?
			\begin{enumerate}
				\item Data drift or newly biased data, as the model's input distribution has shifted under non-stationary conditions.
				\item Server failure, as hardware faults can introduce bias into model outputs through corrupted tensor operations.
				\item Power outage, as interrupted training can introduce bias into model weights through checkpoint corruption.
				\item UI redesign, as changes to the user interface can introduce bias into model inputs through altered data-capture flows.
			\end{enumerate}
		\end{block}
		\begin{exampleblock}{Answer: A}\end{exampleblock}
		\begin{block}{Explanation}
			Bias in outputs typically stems from data or model development. Options B, C, and D are unrelated.
		\end{block}
	\end{frame}

	

	

	


	
	\begin{frame}{Q588 --- Red Teaming}
		\begin{block}{Question}
			What is red teaming in AI?
			\begin{enumerate}
				\item A marketing exercise, as red teaming is used for brand positioning and competitive differentiation.
				\item Adversarial testing by a dedicated team to uncover vulnerabilities and unintended behaviours before deployment.
				\item A financial audit, as red teaming is used for financial controls testing and SOX compliance.
				\item A user survey, as red teaming is used for customer feedback collection and satisfaction measurement.
			\end{enumerate}
		\end{block}
		\begin{exampleblock}{Answer: B}\end{exampleblock}
		\begin{block}{Explanation}
			Red teaming proactively tests for vulnerabilities. Options A, C, and D misstate.
		\end{block}
	\end{frame}

	

	
	% =================================================================
	\section{AI Governance Maturity}
	% =================================================================

	
	\begin{frame}{Q592 --- AI Governance Gaps}
		\begin{block}{Question}
			Which is a common gap in AI governance frameworks?
			\begin{enumerate}
				\item Extensive controls, as frameworks often over-control low-risk use cases and create unnecessary friction.
				\item Missing governance for GenAI, alternative data, or third-party models.
				\item Complete coverage, as frameworks often over-cover all risk dimensions and create unnecessary complexity.
				\item Perfect alignment, as frameworks often over-align with business objectives and create unnecessary constraints.
			\end{enumerate}
		\end{block}
		\begin{exampleblock}{Answer: B}\end{exampleblock}
		\begin{block}{Explanation}
			Common gaps include GenAI, alternative data, and vendor models. Options A, C, and D are not gaps.
		\end{block}
	\end{frame}
	
	\begin{frame}{Q593 --- Governance Roadmap}
		\begin{block}{Question}
			What is the purpose of an AI governance roadmap?
			\begin{enumerate}
				\item A marketing document, as roadmaps are used for brand positioning and competitive differentiation.
				\item A phased plan outlining how the organisation will build, enhance, and sustain AI governance capabilities over time.
				\item A single policy, as roadmaps are policy documents and compliance artefacts.
				\item A code repository, as roadmaps are technical artefacts and infrastructure-as-code modules.
			\end{enumerate}
		\end{block}
		\begin{exampleblock}{Answer: B}\end{exampleblock}
		\begin{block}{Explanation}
			A roadmap is a phased plan. Options A, C, and D misstate.
		\end{block}
	\end{frame}
	
	\begin{frame}{Q594 --- AI Governance KPIs}
		\begin{block}{Question}
			Which is an appropriate KPI for AI governance?
			\begin{enumerate}
				\item Number of employees, as headcount correlates with governance effectiveness and organisational capacity.
				\item Percentage of models validated on schedule.
				\item Office size, as physical footprint correlates with governance effectiveness and operational scale.
				\item Marketing spend, as brand awareness correlates with governance effectiveness and market positioning.
			\end{enumerate}
		\end{block}
		\begin{exampleblock}{Answer: B}\end{exampleblock}
		\begin{block}{Explanation}
			Validation timeliness is a governance KPI. Options A, C, and D are unrelated.
		\end{block}
	\end{frame}

	

	



	
	
	% =================================================================
	\section{Ethics in GenAI}
	% =================================================================
	
	\begin{frame}{Q601 --- GenAI and Autonomy}
		\begin{block}{Question}
			How can GenAI undermine human autonomy?
			\begin{enumerate}
				\item By providing choices, as GenAI expands the option set available to users and increases decision-space dimensionality.
				\item By generating persuasive, tailored content that manipulates users without their awareness.
				\item By explaining decisions, as GenAI provides transparency into decision processes and interpretability artefacts.
				\item By slowing systems, as GenAI introduces latency that reduces user control and increases interaction friction.
			\end{enumerate}
		\end{block}
		\begin{exampleblock}{Answer: B}\end{exampleblock}
		\begin{block}{Explanation}
			Manipulative content undermines autonomy. Options A, C, and D respect or are unrelated to autonomy.
		\end{block}
	\end{frame}
	
	\begin{frame}{Q602 --- GenAI and Bias}
		\begin{block}{Question}
			How does GenAI amplify bias?
			\begin{enumerate}
				\item By generating content that reflects biases present in training data, at scale.
				\item By reducing accuracy, as lower accuracy correlates with higher bias and increased error rates.
				\item By slowing training, as longer training correlates with higher bias and overfitting to minority classes.
				\item By increasing costs, as higher costs correlate with higher bias and increased computational complexity.
			\end{enumerate}
		\end{block}
		\begin{exampleblock}{Answer: A}\end{exampleblock}
		\begin{block}{Explanation}
			GenAI can amplify training-data bias at scale. Options B, C, and D are unrelated.
		\end{block}
	\end{frame}
	

	
	\begin{frame}{Q604 --- GenAI and Consent}
		\begin{block}{Question}
			What consent issues arise with GenAI training data?
			\begin{enumerate}
				\item None, as GenAI training data is always consented and properly licensed.
				\item Data subjects whose content was used in training may not have consented, raising ethical and legal concerns.
				\item Consent is automatically granted, as GenAI training data is publicly available and web-scraped.
				\item Consent is not required, as GenAI training data is exempt from consent requirements under fair-use doctrines.
			\end{enumerate}
		\end{block}
		\begin{exampleblock}{Answer: B}\end{exampleblock}
		\begin{block}{Explanation}
			Training data consent is a critical concern. Options A, C, and D are incorrect.
		\end{block}
	\end{frame}
	
	\begin{frame}{Q605 --- GenAI and Attribution}
		\begin{block}{Question}
			Why is attribution important in GenAI?
			\begin{enumerate}
				\item It isn't important, as GenAI outputs are not attributable and are generated de novo.
				\item To trace sources of content and information, supporting transparency and possibly legal rights such as copyright.
				\item To increase cost, as attribution is primarily a cost optimization technique and FinOps unit-economics lever.
				\item To slow inference, as attribution is primarily a latency optimization technique and inference-time acceleration method.
			\end{enumerate}
		\end{block}
		\begin{exampleblock}{Answer: B}\end{exampleblock}
		\begin{block}{Explanation}
			Attribution supports transparency and copyright compliance. Options A, C, and D are incorrect.
		\end{block}
	\end{frame}
	

	
	\begin{frame}{Q607 --- GenAI and Safety}
		\begin{block}{Question}
			Which is a safety concern with GenAI in critical applications?
			\begin{enumerate}
				\item Faster response, as speed reduces safety in critical applications and increases error propagation.
				\item Generating unsafe or harmful content that could cause physical, psychological, or financial harm.
				\item Lower costs, as cost reduction reduces safety in critical applications and increases technical debt.
				\item Higher accuracy, as accuracy reduces safety in critical applications and increases overconfidence.
			\end{enumerate}
		\end{block}
		\begin{exampleblock}{Answer: B}\end{exampleblock}
		\begin{block}{Explanation}
			Harmful content generation is a safety concern. Options A, C, and D are benefits.
		\end{block}
	\end{frame}
	
	\begin{frame}{Q608 --- GenAI and Consumer Protection}
		\begin{block}{Question}
			Which consumer protection issue arises from GenAI?
			\begin{enumerate}
				\item Increasing prices only, as GenAI directly causes price increases and inflationary pressure.
				\item Misleading or inaccurate information provided to consumers, potentially leading to harm.
				\item Faster service, as speed reduces consumer protection and increases error propagation.
				\item Lower fees, as fee reduction reduces consumer protection and increases hidden costs.
			\end{enumerate}
		\end{block}
		\begin{exampleblock}{Answer: B}\end{exampleblock}
		\begin{block}{Explanation}
			Misleading information is a consumer protection issue. Options A, C, and D are unrelated.
		\end{block}
	\end{frame}
	

	
	\begin{frame}{Q610 --- GenAI and Fraud}
		\begin{block}{Question}
			How can GenAI facilitate fraud?
			\begin{enumerate}
				\item By making it easier to create convincing phishing content and deepfakes at scale.
				\item By improving accuracy, as accuracy reduces fraud and increases detection rates.
				\item By reducing costs, as cost reduction reduces fraud and increases detection efficiency.
				\item By speeding up transactions, as speed reduces fraud and increases detection velocity.
			\end{enumerate}
		\end{block}
		\begin{exampleblock}{Answer: A}\end{exampleblock}
		\begin{block}{Explanation}
			GenAI enables scalable, convincing fraud content. Options B, C, and D are unrelated.
		\end{block}
	\end{frame}
	
	% =================================================================
	\section{Emerging Topics}
	% =================================================================
	

	\begin{frame}{Q616 --- AI and Concentration Risk}
		\begin{block}{Question}
			Why is AI concentration risk a concern?
			\begin{enumerate}
				\item Because it reduces costs, as concentration reduces procurement overhead and increases bargaining power.
				\item Because reliance on a small number of AI providers or models creates correlated failure risk across the ecosystem.
				\item Because it increases accuracy, as concentration improves model quality and increases benchmark performance.
				\item Because it improves speed, as concentration reduces integration complexity and increases deployment velocity.
			\end{enumerate}
		\end{block}
		\begin{exampleblock}{Answer: B}\end{exampleblock}
		\begin{block}{Explanation}
			Concentration amplifies correlated risk. Options A, C, and D are benefits.
		\end{block}
	\end{frame}
	
	
	

	
	\begin{frame}{Q626 --- AI Risk Appetite Cascading}
		\begin{block}{Question}
			How should AI risk appetite cascade through the organisation?
			\begin{enumerate}
				\item It shouldn't, as risk appetite is set independently at each level and is not hierarchical.
				\item From board-set appetite to business-unit tolerances and model-level thresholds, all aligned.
				\item Only at business-unit level, as board appetite is not relevant to AI and is too high-level.
				\item Only at IT level, as AI risk is an IT concern and is managed by technology teams.
			\end{enumerate}
		\end{block}
		\begin{exampleblock}{Answer: B}\end{exampleblock}
		\begin{block}{Explanation}
			Appetite cascades from board to units to models. Options A, C, and D are too narrow.
		\end{block}
	\end{frame}
	
	\begin{frame}{Q627 --- AI Risk Ownership}
		\begin{block}{Question}
			Who owns AI risk in an organisation?
			\begin{enumerate}
				\item The IT department, as AI risk is an IT concern and is managed by technology teams.
				\item Business units that use AI, supported by risk, compliance, and other functions.
				\item Only the CRO, as AI risk is a risk function concern and is managed by risk teams.
				\item Only external vendors, as AI risk is a vendor concern and is managed by third parties.
			\end{enumerate}
		\end{block}
		\begin{exampleblock}{Answer: B}\end{exampleblock}
		\begin{block}{Explanation}
			AI risk ownership sits with business users supported by risk functions. Options A, C, and D are incomplete.
		\end{block}
	\end{frame}
	
	\begin{frame}{Q628 --- AI Escalation Paths}
		\begin{block}{Question}
			What is the purpose of AI escalation paths?
			\begin{enumerate}
				\item To slow down decisions, as escalation introduces delays and bureaucratic friction.
				\item To ensure that AI incidents, breaches, and issues are reported and addressed at the appropriate level promptly.
				\item To reduce transparency, as escalation limits information flow and increases confidentiality.
				\item To bypass governance, as escalation circumvents controls and standard procedures.
			\end{enumerate}
		\end{block}
		\begin{exampleblock}{Answer: B}\end{exampleblock}
		\begin{block}{Explanation}
			Escalation paths ensure timely, appropriate responses. Options A, C, and D are contrary.
		\end{block}
	\end{frame}
	
	\begin{frame}{Q629 --- AI Risk Aggregation}
		\begin{block}{Question}
			What does AI risk aggregation mean?
			\begin{enumerate}
				\item Averaging individual risks, as aggregation is a simple mean calculation and arithmetic average.
				\item Combining individual AI risks into an enterprise-wide view, accounting for interdependencies and correlations.
				\item Ignoring enterprise risk, as aggregation focuses on individual risks and ignores portfolio effects.
				\item Removing individual risks, as aggregation eliminates individual risk and reduces total exposure.
			\end{enumerate}
		\end{block}
		\begin{exampleblock}{Answer: B}\end{exampleblock}
		\begin{block}{Explanation}
			Aggregation considers interdependencies. Options A, C, and D misstate.
		\end{block}
	\end{frame}
	
	\begin{frame}{Q630 --- AI Governance Success Metrics}
		\begin{block}{Question}
			Which is a good metric for AI governance success?
			\begin{enumerate}
				\item Number of meetings, as meeting frequency correlates with governance effectiveness and organisational alignment.
				\item Reduction in AI incidents, improvements in validation timeliness, and increasing compliance with governance standards.
				\item Number of documents, as documentation volume correlates with governance effectiveness and process maturity.
				\item Cost of IT, as IT spend correlates with governance effectiveness and operational efficiency.
			\end{enumerate}
		\end{block}
		\begin{exampleblock}{Answer: B}\end{exampleblock}
		\begin{block}{Explanation}
			Governance success metrics track incidents, timeliness, and compliance. Options A, C, and D are not meaningful metrics.
		\end{block}
	\end{frame}
	
	% =================================================================
	\section{Ethics and Regulation in Depth}
	% =================================================================

	
	
\end{document}